\documentclass[11pt,a4paper]{article}
\usepackage[utf8]{inputenc}
\usepackage[T1]{fontenc}
\usepackage[spanish,es-minimal,es-noshorthands]{babel}
\usepackage{lmodern}
\usepackage[sfdefault,lining]{FiraSans}
\usepackage[a4paper,top=2.3cm,bottom=2.2cm,left=1.8cm,right=1.8cm,headheight=1.3cm,headsep=0.4cm,footskip=1cm]{geometry}
\usepackage{amsmath,amssymb}
\usepackage{icomma}
\usepackage{graphicx}
\usepackage[table]{xcolor}
\usepackage{tikz}
\usepackage{fontawesome5}
\usepackage{ccicons}
\usepackage{fancyhdr}
\usepackage{lastpage}
\usepackage{tasks}
\usepackage{needspace}
\usepackage[most]{tcolorbox}
\usepackage{array}
\usepackage[hidelinks]{hyperref}
\AtBeginDocument{\hypersetup{pdftitle={Soluciones de la ficha de la unidad 4: Fracciones. 1º ESO},pdfsubject={Matemáticas 1º ESO},pdfkeywords={fracciones, fracción de un número, porcentajes, fracciones equivalentes, común denominador, comparar fracciones, operaciones con fracciones, problemas}}}

\definecolor{navy}{HTML}{1B2748}
\definecolor{accent}{HTML}{E8603A}
\definecolor{teal}{HTML}{1F9E90}
\definecolor{soft}{HTML}{EEF1F8}
\definecolor{muted}{HTML}{5D6685}

\setlength{\parindent}{0pt}
\setlength{\parskip}{3pt}
\settasks{label=\alph*), label-format=\bfseries\color{navy}, label-width=1.3em, item-indent=1.8em, column-sep=0.6em, after-item-skip=5pt, before-skip=3pt, after-skip=4pt}

% ---------- capçalera i peu ----------
\newcommand{\llicencia}{\href{https://creativecommons.org/licenses/by-nc-sa/4.0/deed.es}{\ccbyncsa}\ \ Licencia Creative Commons Atribución-NoComercial-CompartirIgual 4.0 Internacional (CC BY-NC-SA 4.0)}
\pagestyle{fancy}
\fancyhf{}
\fancyhead[L]{\raisebox{-0.35\height}{\includegraphics[height=0.95cm]{logo-gva.png}}}
\fancyhead[C]{\small\color{muted}Soluciones de la ficha · Unidad 4: Fracciones · 1\textsuperscript{o} ESO}
\fancyhead[R]{\raisebox{-0.35\height}{\includegraphics[height=0.95cm]{logo-vermellar.png}}}
\fancyfoot[L]{\scriptsize\color{muted}\llicencia}
\fancyfoot[R]{\small\color{muted}\thepage}
\renewcommand{\headrulewidth}{0.4pt}
\renewcommand{\headrule}{\color{muted!40}\hrule height \headrulewidth}
\fancypagestyle{primera}{\fancyhf{}\renewcommand{\headrulewidth}{0pt}
  \fancyfoot[L]{\scriptsize\color{muted}\llicencia}
  \fancyfoot[R]{\small\color{muted}\thepage}}

% ---------- nivells i exercicis ----------
\newcommand{\estrelles}[1]{{\color{accent}\ifcase#1\or\faStar\faStar[regular]\faStar[regular]\or\faStar\faStar\faStar[regular]\or\faStar\faStar\faStar\fi}}
\newcounter{exer}
\newcommand{\ex}[1]{\par\needspace{4\baselineskip}\medskip\stepcounter{exer}\noindent{\color{navy}\bfseries\theexer.}\ \ {\small\estrelles{#1}}\ \ }

\newcommand{\seccio}[2]{\par\needspace{10\baselineskip}\bigskip
  \begin{tcolorbox}[colback=navy,colframe=navy,coltext=white,arc=2mm,boxrule=0pt,left=3mm,right=3mm,top=1.2mm,bottom=1.2mm]
  {\large\bfseries #1\quad #2}
  \end{tcolorbox}\vspace{-2pt}}

\newcolumntype{C}[1]{>{\centering\arraybackslash}p{#1}}

% quadre de recordatori (només a la fitxa d'exercicis)
\newcommand{\recorda}[1]{\par\smallskip\begin{tcolorbox}[colback=accent!7,colframe=accent!55,arc=2mm,boxrule=0.5pt,left=2.5mm,right=2.5mm,top=1mm,bottom=1mm]\small{\color{accent}\faLightbulb[regular]}\ \textbf{Recuerda:}\ #1\end{tcolorbox}}
\newcommand{\res}[1]{\boxed{#1}}

% =====================================================================
% Capçalera FentMates de les fitxes: un color per a cada curs.
% Els colors són els de la pàgina inicial (LEVELS a assets/js/materials.js);
% el logo ve de _eines/logo/pdf/ (python3 _eines/logo/genera.py).
%
% Ús, al final del preàmbul de la unitat:
%   % =====================================================================
% Capçalera FentMates de les fitxes: un color per a cada curs.
% Els colors són els de la pàgina inicial (LEVELS a assets/js/materials.js);
% el logo ve de _eines/logo/pdf/ (python3 _eines/logo/genera.py).
%
% Ús, al final del preàmbul de la unitat:
%   % =====================================================================
% Capçalera FentMates de les fitxes: un color per a cada curs.
% Els colors són els de la pàgina inicial (LEVELS a assets/js/materials.js);
% el logo ve de _eines/logo/pdf/ (python3 _eines/logo/genera.py).
%
% Ús, al final del preàmbul de la unitat:
%   \input{../comu/fentmates.tex}
%   \fmcurs{1eso}
% i, al començament del document:
%   \capcalera{Matemàtiques}{1r ESO}{Fitxa d'exercicis · Unitat 2}{Divisibilitat}
% Opcional abans de \input: \def\fmneed{14} (espai mínim abans de cada secció).
% =====================================================================
\makeatletter
\def\fm@logodir{../../logo/pdf}
\providecommand{\fmneed}{10}
\def\fm@cursnom{}

% clau del curs, color i color del text damunt del color
\newcommand{\fm@curs}[3]{\@namedef{fm@color@#1}{#2}\@namedef{fm@text@#1}{#3}}
\fm@curs{1eso}{E6B84A}{navy}
\fm@curs{2eso}{3B82F6}{white}
\fm@curs{3eso}{FF7A50}{white}
\fm@curs{4eso-a}{A78BFA}{white}
\fm@curs{4eso-b}{6D28D9}{white}
\fm@curs{1batx-mi}{9B1C1C}{white}
\fm@curs{1batx-ccss}{EF6B6B}{white}
\fm@curs{2batx-mii}{1E6B47}{white}
\fm@curs{2batx-ccss}{2BB3A3}{white}

\newcommand{\fmcurs}[1]{%
  \@ifundefined{fm@color@#1}{\PackageError{fentmates}{Curs desconegut: #1}{}}{}%
  \edef\fm@tmp{\noexpand\definecolor{curs}{HTML}{\@nameuse{fm@color@#1}}%
               \noexpand\colorlet{curstext}{\@nameuse{fm@text@#1}}}\fm@tmp
  \colorlet{cursfosc}{curs!72!black}%
  \def\fm@logo{\fm@logodir/vertical-#1.pdf}%
  \def\fm@marca{\fm@logodir/marca-#1.pdf}}

% ---------- capçalera de la primera pàgina ----------
% #1 matèria (petit), #2 curs, #3 tipus de document i unitat, #4 títol
\newcommand{\capcalera}[4]{%
  \gdef\fm@cursnom{#2}%
  \par\noindent
  \begin{tcolorbox}[colback=curs,colframe=curs,boxrule=0pt,arc=2.2mm,width=3.5cm,height=2.15cm,
      valign=center,left=3mm,right=2mm,top=1mm,bottom=1mm,nobeforeafter,box align=center]
    \color{curstext}{\footnotesize #1}\par\vspace{1pt}{\fontsize{19}{21}\selectfont\bfseries #2\par}%
  \end{tcolorbox}\hspace{5mm}%
  \begin{minipage}[c]{\dimexpr\linewidth-3.5cm-5mm-2.9cm\relax}
    {\color{cursfosc}\small\bfseries\MakeUppercase{#3}}\par\vspace{2pt}
    {\color{navy}\fontsize{20}{23}\selectfont\bfseries #4\par}
  \end{minipage}\hfill
  \begin{minipage}[c]{2.6cm}\raggedleft\includegraphics[width=2.6cm]{\fm@logo}\end{minipage}\par
  \vspace{2pt}}

% ---------- barres de secció del color del curs ----------
\renewcommand{\seccio}[2]{\par\needspace{\fmneed\baselineskip}\bigskip
  \begin{tcolorbox}[colback=curs,colframe=curs,coltext=curstext,arc=2mm,boxrule=0pt,left=3mm,right=3mm,top=1.2mm,bottom=1.2mm]
  {\large\bfseries #1\quad #2}
  \end{tcolorbox}\vspace{-2pt}}

% ---------- estrelles de dificultat del color del curs (a Batxillerat no se'n posen) ----------
\@ifundefined{estrelles}{}{%
  \renewcommand{\estrelles}[1]{{\color{curs}\ifcase#1\or\faStar\faStar[regular]\faStar[regular]\or\faStar\faStar\faStar[regular]\or\faStar\faStar\faStar\fi}}}

% ---------- capçalera i peu de les pàgines ----------
% primera pàgina: logos institucionals al peu
\fancypagestyle{primera}{\fancyhf{}\renewcommand{\headrulewidth}{0pt}
  \fancyfoot[L]{\raisebox{-0.32\height}{\includegraphics[height=0.78cm]{logo-gva.png}}\hspace{5mm}%
    \raisebox{-0.32\height}{\includegraphics[height=0.78cm]{logo-vermellar.png}}}
  \fancyfoot[R]{\scriptsize\color{muted}\href{https://creativecommons.org/licenses/by-nc-sa/4.0/deed.ca}{\ccbyncsa}\ \ CC BY-NC-SA 4.0 · fentmates.org\qquad\small\thepage}}
% la resta: el símbol a l'esquerra i el curs a la dreta, amb el títol al centre com abans
\fancyhead[L]{\raisebox{-0.3\height}{\includegraphics[height=0.62cm]{\fm@marca}}}
\fancyhead[R]{\small\bfseries\color{cursfosc}\fm@cursnom}
\renewcommand{\headrule}{\color{curs!45}\hrule height \headrulewidth}
\makeatother

%   \fmcurs{1eso}
% i, al començament del document:
%   \capcalera{Matemàtiques}{1r ESO}{Fitxa d'exercicis · Unitat 2}{Divisibilitat}
% Opcional abans de \input: \def\fmneed{14} (espai mínim abans de cada secció).
% =====================================================================
\makeatletter
\def\fm@logodir{../../logo/pdf}
\providecommand{\fmneed}{10}
\def\fm@cursnom{}

% clau del curs, color i color del text damunt del color
\newcommand{\fm@curs}[3]{\@namedef{fm@color@#1}{#2}\@namedef{fm@text@#1}{#3}}
\fm@curs{1eso}{E6B84A}{navy}
\fm@curs{2eso}{3B82F6}{white}
\fm@curs{3eso}{FF7A50}{white}
\fm@curs{4eso-a}{A78BFA}{white}
\fm@curs{4eso-b}{6D28D9}{white}
\fm@curs{1batx-mi}{9B1C1C}{white}
\fm@curs{1batx-ccss}{EF6B6B}{white}
\fm@curs{2batx-mii}{1E6B47}{white}
\fm@curs{2batx-ccss}{2BB3A3}{white}

\newcommand{\fmcurs}[1]{%
  \@ifundefined{fm@color@#1}{\PackageError{fentmates}{Curs desconegut: #1}{}}{}%
  \edef\fm@tmp{\noexpand\definecolor{curs}{HTML}{\@nameuse{fm@color@#1}}%
               \noexpand\colorlet{curstext}{\@nameuse{fm@text@#1}}}\fm@tmp
  \colorlet{cursfosc}{curs!72!black}%
  \def\fm@logo{\fm@logodir/vertical-#1.pdf}%
  \def\fm@marca{\fm@logodir/marca-#1.pdf}}

% ---------- capçalera de la primera pàgina ----------
% #1 matèria (petit), #2 curs, #3 tipus de document i unitat, #4 títol
\newcommand{\capcalera}[4]{%
  \gdef\fm@cursnom{#2}%
  \par\noindent
  \begin{tcolorbox}[colback=curs,colframe=curs,boxrule=0pt,arc=2.2mm,width=3.5cm,height=2.15cm,
      valign=center,left=3mm,right=2mm,top=1mm,bottom=1mm,nobeforeafter,box align=center]
    \color{curstext}{\footnotesize #1}\par\vspace{1pt}{\fontsize{19}{21}\selectfont\bfseries #2\par}%
  \end{tcolorbox}\hspace{5mm}%
  \begin{minipage}[c]{\dimexpr\linewidth-3.5cm-5mm-2.9cm\relax}
    {\color{cursfosc}\small\bfseries\MakeUppercase{#3}}\par\vspace{2pt}
    {\color{navy}\fontsize{20}{23}\selectfont\bfseries #4\par}
  \end{minipage}\hfill
  \begin{minipage}[c]{2.6cm}\raggedleft\includegraphics[width=2.6cm]{\fm@logo}\end{minipage}\par
  \vspace{2pt}}

% ---------- barres de secció del color del curs ----------
\renewcommand{\seccio}[2]{\par\needspace{\fmneed\baselineskip}\bigskip
  \begin{tcolorbox}[colback=curs,colframe=curs,coltext=curstext,arc=2mm,boxrule=0pt,left=3mm,right=3mm,top=1.2mm,bottom=1.2mm]
  {\large\bfseries #1\quad #2}
  \end{tcolorbox}\vspace{-2pt}}

% ---------- estrelles de dificultat del color del curs (a Batxillerat no se'n posen) ----------
\@ifundefined{estrelles}{}{%
  \renewcommand{\estrelles}[1]{{\color{curs}\ifcase#1\or\faStar\faStar[regular]\faStar[regular]\or\faStar\faStar\faStar[regular]\or\faStar\faStar\faStar\fi}}}

% ---------- capçalera i peu de les pàgines ----------
% primera pàgina: logos institucionals al peu
\fancypagestyle{primera}{\fancyhf{}\renewcommand{\headrulewidth}{0pt}
  \fancyfoot[L]{\raisebox{-0.32\height}{\includegraphics[height=0.78cm]{logo-gva.png}}\hspace{5mm}%
    \raisebox{-0.32\height}{\includegraphics[height=0.78cm]{logo-vermellar.png}}}
  \fancyfoot[R]{\scriptsize\color{muted}\href{https://creativecommons.org/licenses/by-nc-sa/4.0/deed.ca}{\ccbyncsa}\ \ CC BY-NC-SA 4.0 · fentmates.org\qquad\small\thepage}}
% la resta: el símbol a l'esquerra i el curs a la dreta, amb el títol al centre com abans
\fancyhead[L]{\raisebox{-0.3\height}{\includegraphics[height=0.62cm]{\fm@marca}}}
\fancyhead[R]{\small\bfseries\color{cursfosc}\fm@cursnom}
\renewcommand{\headrule}{\color{curs!45}\hrule height \headrulewidth}
\makeatother

%   \fmcurs{1eso}
% i, al començament del document:
%   \capcalera{Matemàtiques}{1r ESO}{Fitxa d'exercicis · Unitat 2}{Divisibilitat}
% Opcional abans de \input: \def\fmneed{14} (espai mínim abans de cada secció).
% =====================================================================
\makeatletter
\def\fm@logodir{../../logo/pdf}
\providecommand{\fmneed}{10}
\def\fm@cursnom{}

% clau del curs, color i color del text damunt del color
\newcommand{\fm@curs}[3]{\@namedef{fm@color@#1}{#2}\@namedef{fm@text@#1}{#3}}
\fm@curs{1eso}{E6B84A}{navy}
\fm@curs{2eso}{3B82F6}{white}
\fm@curs{3eso}{FF7A50}{white}
\fm@curs{4eso-a}{A78BFA}{white}
\fm@curs{4eso-b}{6D28D9}{white}
\fm@curs{1batx-mi}{9B1C1C}{white}
\fm@curs{1batx-ccss}{EF6B6B}{white}
\fm@curs{2batx-mii}{1E6B47}{white}
\fm@curs{2batx-ccss}{2BB3A3}{white}

\newcommand{\fmcurs}[1]{%
  \@ifundefined{fm@color@#1}{\PackageError{fentmates}{Curs desconegut: #1}{}}{}%
  \edef\fm@tmp{\noexpand\definecolor{curs}{HTML}{\@nameuse{fm@color@#1}}%
               \noexpand\colorlet{curstext}{\@nameuse{fm@text@#1}}}\fm@tmp
  \colorlet{cursfosc}{curs!72!black}%
  \def\fm@logo{\fm@logodir/vertical-#1.pdf}%
  \def\fm@marca{\fm@logodir/marca-#1.pdf}}

% ---------- capçalera de la primera pàgina ----------
% #1 matèria (petit), #2 curs, #3 tipus de document i unitat, #4 títol
\newcommand{\capcalera}[4]{%
  \gdef\fm@cursnom{#2}%
  \par\noindent
  \begin{tcolorbox}[colback=curs,colframe=curs,boxrule=0pt,arc=2.2mm,width=3.5cm,height=2.15cm,
      valign=center,left=3mm,right=2mm,top=1mm,bottom=1mm,nobeforeafter,box align=center]
    \color{curstext}{\footnotesize #1}\par\vspace{1pt}{\fontsize{19}{21}\selectfont\bfseries #2\par}%
  \end{tcolorbox}\hspace{5mm}%
  \begin{minipage}[c]{\dimexpr\linewidth-3.5cm-5mm-2.9cm\relax}
    {\color{cursfosc}\small\bfseries\MakeUppercase{#3}}\par\vspace{2pt}
    {\color{navy}\fontsize{20}{23}\selectfont\bfseries #4\par}
  \end{minipage}\hfill
  \begin{minipage}[c]{2.6cm}\raggedleft\includegraphics[width=2.6cm]{\fm@logo}\end{minipage}\par
  \vspace{2pt}}

% ---------- barres de secció del color del curs ----------
\renewcommand{\seccio}[2]{\par\needspace{\fmneed\baselineskip}\bigskip
  \begin{tcolorbox}[colback=curs,colframe=curs,coltext=curstext,arc=2mm,boxrule=0pt,left=3mm,right=3mm,top=1.2mm,bottom=1.2mm]
  {\large\bfseries #1\quad #2}
  \end{tcolorbox}\vspace{-2pt}}

% ---------- estrelles de dificultat del color del curs (a Batxillerat no se'n posen) ----------
\@ifundefined{estrelles}{}{%
  \renewcommand{\estrelles}[1]{{\color{curs}\ifcase#1\or\faStar\faStar[regular]\faStar[regular]\or\faStar\faStar\faStar[regular]\or\faStar\faStar\faStar\fi}}}

% ---------- capçalera i peu de les pàgines ----------
% primera pàgina: logos institucionals al peu
\fancypagestyle{primera}{\fancyhf{}\renewcommand{\headrulewidth}{0pt}
  \fancyfoot[L]{\raisebox{-0.32\height}{\includegraphics[height=0.78cm]{logo-gva.png}}\hspace{5mm}%
    \raisebox{-0.32\height}{\includegraphics[height=0.78cm]{logo-vermellar.png}}}
  \fancyfoot[R]{\scriptsize\color{muted}\href{https://creativecommons.org/licenses/by-nc-sa/4.0/deed.ca}{\ccbyncsa}\ \ CC BY-NC-SA 4.0 · fentmates.org\qquad\small\thepage}}
% la resta: el símbol a l'esquerra i el curs a la dreta, amb el títol al centre com abans
\fancyhead[L]{\raisebox{-0.3\height}{\includegraphics[height=0.62cm]{\fm@marca}}}
\fancyhead[R]{\small\bfseries\color{cursfosc}\fm@cursnom}
\renewcommand{\headrule}{\color{curs!45}\hrule height \headrulewidth}
\makeatother

\fmcurs{1eso}

\begin{document}
\thispagestyle{primera}

\capcalera{Matemáticas}{1\textsuperscript{o} ESO}{Soluciones de la ficha · Unidad 4}{Fracciones}
\vspace{8pt}

\begin{tcolorbox}[colback=soft,colframe=soft,arc=2mm,left=3mm,right=3mm,top=1.5mm,bottom=1.5mm]
\small Intenta resolver cada ejercicio antes de mirar su solución. Si no te sale el mismo resultado, repasa los pasos o consulta los apuntes interactivos de la unidad. Recuerda que los resultados se tienen que dar siempre en forma de fracción irreducible.
\end{tcolorbox}

% =====================================================
\seccio{1}{El concepto de fracción}

\ex{1} \begin{tasks}(6)
  \task $\dfrac{4}{10}$
  \task $\dfrac{3}{4}$
  \task $\dfrac{1}{7}$
  \task $\dfrac{2}{4}$
  \task $\dfrac{5}{7}$
  \task $\dfrac{4}{6}$
\end{tasks}

\ex{1} \begin{tasks}(2)
  \task num.\ 3, den.\ 5: tres quintos
  \task num.\ 4, den.\ 6: cuatro sextos
  \task num.\ 12, den.\ 15: doce quinceavos
  \task num.\ 4, den.\ 3: cuatro tercios
\end{tasks}

\ex{1} \begin{tasks}(3)
  \task $\dfrac{3}{5}$\ \ \begin{tikzpicture}[baseline=(current bounding box.center)]\fill[accent!75] (0.000,0) rectangle ++(0.700,0.55);\fill[accent!75] (0.700,0) rectangle ++(0.700,0.55);\fill[accent!75] (1.400,0) rectangle ++(0.700,0.55);\draw (0.700,0)--(0.700,0.55);\draw (1.400,0)--(1.400,0.55);\draw (2.100,0)--(2.100,0.55);\draw (2.800,0)--(2.800,0.55);\draw[thick] (0,0) rectangle (3.5,0.55);\end{tikzpicture}
  \task $\dfrac{5}{8}$\ \ \begin{tikzpicture}[scale=0.42,baseline=(current bounding box.center)]\fill[accent!75] (0,1) rectangle ++(1,1);\fill[accent!75] (1,1) rectangle ++(1,1);\fill[accent!75] (2,1) rectangle ++(1,1);\fill[accent!75] (3,1) rectangle ++(1,1);\fill[accent!75] (0,0) rectangle ++(1,1);\draw[thick] (0,0) grid (4,2);\end{tikzpicture}
  \task $\dfrac{2}{3}$\ \ \begin{tikzpicture}[baseline=(current bounding box.center)]\fill[accent!75] (0,0) -- (90.00:0.65) arc[start angle=90.00,end angle=-30.00,radius=0.65] -- cycle;\fill[accent!75] (0,0) -- (-30.00:0.65) arc[start angle=-30.00,end angle=-150.00,radius=0.65] -- cycle;\draw[thick] (0,0) circle (0.65);\draw[thick] (0,0) -- (90.00:0.65);\draw[thick] (0,0) -- (-30.00:0.65);\draw[thick] (0,0) -- (-150.00:0.65);\end{tikzpicture}
\end{tasks}

\ex{1} Dividimos cada unidad en 10 partes para $\dfrac{6}{10}$ y en 4 para $\dfrac{7}{4}=1+\dfrac{3}{4}$:\par\smallskip{\centering \begin{tikzpicture}[baseline=0]\draw[-latex,thick] (-0.3,0)--(9.50,0);\draw[muted] (1.125,0.1)--(1.125,-0.1);\draw[muted] (2.250,0.1)--(2.250,-0.1);\draw[muted] (3.375,0.1)--(3.375,-0.1);\draw[muted] (5.625,0.1)--(5.625,-0.1);\draw[muted] (6.750,0.1)--(6.750,-0.1);\draw[muted] (7.875,0.1)--(7.875,-0.1);\draw[thick] (0.00,0.18)--(0.00,-0.18) node[below=1pt,font=\small]{$0$};\draw[thick] (4.50,0.18)--(4.50,-0.18) node[below=1pt,font=\small]{$1$};\draw[thick] (9.00,0.18)--(9.00,-0.18) node[below=1pt,font=\small]{$2$};\fill[accent] (2.700,0) circle (0.09) node[above=3pt,text=accent]{$\frac{6}{10}$};\fill[accent] (7.875,0) circle (0.09) node[below=10pt=3pt,text=accent]{$\frac{7}{4}$};\end{tikzpicture}\par}

\ex{2} \begin{tasks}(3)
  \task Impropia: $2+\dfrac{1}{3}$
  \task Propia
  \task Igual a 1
  \task Impropia: $3+\dfrac{1}{4}$
  \task Propia
  \task Impropia: $3+\dfrac{2}{6}$
\end{tasks}

\ex{2} \begin{tasks}(4)
  \task $\dfrac{15}{60}$
  \task $\dfrac{4}{12}$
  \task $\dfrac{5}{7}$
  \task $\dfrac{300}{1000}$
\end{tasks}

\ex{3} Entre $\dfrac{9}{9}$ y $\dfrac{18}{9}$: \ $\res{\dfrac{10}{9},\ \dfrac{11}{9},\ \dfrac{12}{9},\ \dfrac{13}{9},\ \dfrac{14}{9},\ \dfrac{15}{9},\ \dfrac{16}{9},\ \dfrac{17}{9}}$. \ Entre $\dfrac{12}{6}$ y $\dfrac{18}{6}$: \ $\res{\dfrac{13}{6},\ \dfrac{14}{6},\ \dfrac{15}{6},\ \dfrac{16}{6},\ \dfrac{17}{6}}$.


% =====================================================
\seccio{2}{Fracción de un número}

\ex{1} \begin{tasks}(3)
  \task $(48:6)\cdot5=\res{40}$
  \task $(490:7)\cdot4=\res{280}$
  \task $(200:5)\cdot3=\res{120}$
  \task $(32:8)\cdot7=\res{28}$
  \task $(60:5)\cdot2=\res{24}$
  \task $(180:6)\cdot3=\res{90}$
\end{tasks}

\ex{1} Chicos: $(30:5)\cdot2=\res{12}$. \ Chicas: $30-12=\res{18}$.

\ex{2} $(2400:8)\cdot3=\res{900\text{ L}}$. \ Faltan $2400-900=\res{1500\text{ L}}$.

\ex{2} \begin{tasks}(2)
  \task $36:3=12$; \ $12\cdot4=\res{48}$
  \task $18:2=9$; \ $9\cdot7=\res{63}$
\end{tasks}

\ex{2} Le quedan $\dfrac{2}{5}$ de la paga, que son 14\,€: \ $14:2=7$; \ $7\cdot5=\res{35\text{\,€}}$.

\ex{3} $\dfrac{3}{4}$ de $360=270$; \ $\dfrac{2}{3}$ de $270=\res{180}$.

\ex{3} 1\textsuperscript{o}: $(84:7)\cdot3=\res{36}$. \ 2\textsuperscript{o}: $84:4=\res{21}$. \ 3\textsuperscript{o}: $84-36-21=\res{27}$.


% =====================================================
\seccio{3}{Porcentajes y fracciones}

\ex{1} \begin{tasks}(3)
  \task $\dfrac{20\cdot180}{100}=\res{36}$
  \task $\dfrac{45\cdot80}{100}=\res{36}$
  \task $\dfrac{12\cdot350}{100}=\res{42}$
  \task $\dfrac{85\cdot1300}{100}=\res{1105}$
  \task $\dfrac{50\cdot64}{100}=\res{32}$
  \task $\dfrac{25\cdot120}{100}=\res{30}$
\end{tasks}

\ex{1} \begin{tasks}(4)
  \task $\dfrac{50}{100}=\dfrac{1}{2}$
  \task $\dfrac{25}{100}=\dfrac{1}{4}$
  \task $\dfrac{75}{100}=\dfrac{3}{4}$
  \task $\dfrac{20}{100}=\dfrac{1}{5}$
  \task $\dfrac{10}{100}=\dfrac{1}{10}$
  \task $\dfrac{40}{100}=\dfrac{2}{5}$
  \task $\dfrac{5}{100}=\dfrac{1}{20}$
  \task $\dfrac{35}{100}=\dfrac{7}{20}$
\end{tasks}

\ex{1} \begin{tasks}(3)
  \task $\dfrac{1}{2}=\dfrac{50}{100}=\res{50\,\%}$
  \task $\dfrac{3}{4}=\dfrac{75}{100}=\res{75\,\%}$
  \task $\dfrac{2}{5}=\dfrac{40}{100}=\res{40\,\%}$
  \task $\dfrac{7}{10}=\dfrac{70}{100}=\res{70\,\%}$
  \task $\dfrac{9}{20}=\dfrac{45}{100}=\res{45\,\%}$
  \task $\dfrac{3}{25}=\dfrac{12}{100}=\res{12\,\%}$
\end{tasks}

\ex{2} $25\,\%=\dfrac{1}{4}$: \ descuento $80:4=\res{20\text{\,€}}$; \ pagamos $80-20=\res{60\text{\,€}}$.

\ex{2} Chicas: $\dfrac{52\cdot650}{100}=\res{338}$. \ Chicos: $650-338=\res{312}$.

\ex{3} De aquí: $100-7-25-12=\res{56\,\%}$. \ Personas: $\res{26\,460}$, \ $\res{94\,500}$, \ $\res{45\,360}$ \ y \ $\res{211\,680}$.

\ex{3} Han aprobado $17$: \ $\dfrac{17}{25}=\dfrac{68}{100}=\res{68\,\%}$.


% =====================================================
\seccio{4}{Fracciones equivalentes}

\ex{1} \begin{tasks}(2)
  \task $\dfrac{1}{4}=\dfrac{2}{8}=\dfrac{3}{12}=\dfrac{4}{16}=\dfrac{5}{20}$
  \task $\dfrac{2}{5}=\dfrac{4}{10}=\dfrac{6}{15}=\dfrac{8}{20}=\dfrac{10}{25}$
  \task $\dfrac{5}{15}=\dfrac{10}{30}=\dfrac{15}{45}=\dfrac{20}{60}=\dfrac{25}{75}$
  \task $\dfrac{7}{13}=\dfrac{14}{26}=\dfrac{21}{39}=\dfrac{28}{52}=\dfrac{35}{65}$
  \task $\dfrac{11}{17}=\dfrac{22}{34}=\dfrac{33}{51}=\dfrac{44}{68}=\dfrac{55}{85}$
  \task $\dfrac{1}{10}=\dfrac{2}{20}=\dfrac{3}{30}=\dfrac{4}{40}=\dfrac{5}{50}$
\end{tasks}

\ex{1} \begin{tasks}(3)
  \task $\dfrac{4}{12}=\dfrac{2}{6}=\dfrac{1}{3}$
  \task $\dfrac{12}{20}=\dfrac{6}{10}=\dfrac{3}{5}$
  \task $\dfrac{66}{84}=\dfrac{33}{42}=\dfrac{11}{14}$
  \task $\dfrac{15}{45}=\dfrac{5}{15}=\dfrac{1}{3}$
  \task $\dfrac{50}{100}=\dfrac{25}{50}=\dfrac{1}{2}$
\end{tasks}\par (Hay otras respuestas posibles.)

\ex{1} \begin{tasks}(3)
  \task No: $\dfrac{5}{7}$
  \task Sí: $\text{m.c.d.}(5,7)=1$
  \task No: $\dfrac{2}{1}=2$
\end{tasks}

\ex{2} \begin{tasks}(5)
  \task $\dfrac{4}{6}=\res{\dfrac{2}{3}}$
  \task $\dfrac{4}{20}=\res{\dfrac{1}{5}}$
  \task $\dfrac{3}{12}=\res{\dfrac{1}{4}}$
  \task $\dfrac{2}{6}=\res{\dfrac{1}{3}}$
  \task $\dfrac{35}{40}=\res{\dfrac{7}{8}}$
  \task $\dfrac{75}{100}=\res{\dfrac{3}{4}}$
  \task $\dfrac{4}{16}=\res{\dfrac{1}{4}}$
  \task $\dfrac{14}{21}=\res{\dfrac{2}{3}}$
  \task $\dfrac{10}{15}=\res{\dfrac{2}{3}}$
  \task $\dfrac{25}{45}=\res{\dfrac{5}{9}}$
  \task $\dfrac{13}{52}=\res{\dfrac{1}{4}}$
  \task $\dfrac{33}{34}=\res{\dfrac{33}{34}}$
  \task $\dfrac{705}{120}=\res{\dfrac{47}{8}}$
  \task $\dfrac{170}{105}=\res{\dfrac{34}{21}}$
  \task $\dfrac{256}{512}=\res{\dfrac{1}{2}}$
\end{tasks}

\ex{2} \begin{tasks}(3)
  \task $\dfrac{54}{24}=\dfrac{9}{4}$
  \task $\dfrac{40}{88}=\dfrac{5}{11}$
  \task $\dfrac{81}{36}=\dfrac{9}{4}$
  \task $\dfrac{72}{32}=\dfrac{9}{4}$
  \task $\dfrac{75}{185}=\dfrac{15}{37}$
  \task $\dfrac{50}{110}=\dfrac{5}{11}$
\end{tasks}\par Equivalentes: $\dfrac{54}{24}$, $\dfrac{81}{36}$ y $\dfrac{72}{32}$ (todas valen $\dfrac{9}{4}$); \ $\dfrac{40}{88}$ y $\dfrac{50}{110}$ (valen $\dfrac{5}{11}$).

\ex{3} $\dfrac{15}{24}=\dfrac{5}{8}=\res{\dfrac{10}{16}}$ \qquad $\dfrac{5}{8}=\dfrac{5\cdot16}{8\cdot16}=\res{\dfrac{80}{128}}$

\ex{3} $\res{\dfrac{14}{21},\ \dfrac{16}{24},\ \dfrac{18}{27},\ \dfrac{20}{30}}$


% =====================================================
\seccio{5}{Comprobar si dos fracciones son equivalentes}

\ex{1} \begin{tasks}(2)
  \task $8\cdot12=96$, \ $32\cdot3=96$: \ \res{\text{sí}}
  \task $6\cdot34=204$, \ $15\cdot26=390$: \ \res{\text{no}}
  \task $7\cdot46=322$, \ $23\cdot14=322$: \ \res{\text{sí}}
  \task $1\cdot30=30$, \ $2\cdot15=30$: \ \res{\text{sí}}
  \task $68\cdot13=884$, \ $52\cdot17=884$: \ \res{\text{sí}}
  \task $3\cdot67=201$, \ $29\cdot33=957$: \ \res{\text{no}}
\end{tasks}

\ex{1} \begin{tasks}(5)
  \task $x=\res{9}$
  \task $x=\res{32}$
  \task $x=\res{3}$
  \task $x=\res{98}$
  \task $x=\res{33}$
\end{tasks}

\ex{2} $\dfrac{2}{5}=\dfrac{\res{8}}{20}=\dfrac{12}{\res{30}}=\dfrac{\res{18}}{45}$

\ex{2} $\res{\dfrac{14}{20}}$ ($=\dfrac{7}{10}$); \ las demás valen $\dfrac{2}{3}$.

\ex{3} $\dfrac{15}{20}=\dfrac{3}{4}$: \ $x=\res{9}$, \ $y=\res{12}$.

\ex{3} En $\dfrac{3}{5}$ la diferencia es $2$; amplificamos por $24:2=12$: \ $\res{\dfrac{36}{60}}$.


% =====================================================
\seccio{6}{Reducción a común denominador}

\ex{1} \begin{tasks}(2)
  \task $\dfrac{7}{35},\ \dfrac{15}{35}$
  \task $\dfrac{4}{18},\ \dfrac{3}{18}$
  \task $\dfrac{6}{8},\ \dfrac{3}{8}$
  \task $\dfrac{8}{28},\ \dfrac{7}{28}$
  \task $\dfrac{10}{42},\ \dfrac{63}{42}$
  \task $\dfrac{6}{10},\ \dfrac{1}{10}$
  \task $\dfrac{3}{12},\ \dfrac{2}{12}$
  \task $\dfrac{3}{7},\ \dfrac{35}{7}$
\end{tasks}

\ex{2} \begin{tasks}(2)
  \task $\text{m.c.m.}=48$: \ $\dfrac{14}{48},\ \dfrac{9}{48}$
  \task $\text{m.c.m.}=24$: \ $\dfrac{12}{24},\ \dfrac{16}{24},\ \dfrac{15}{24}$
  \task $\text{m.c.m.}=30$: \ $\dfrac{105}{30},\ \dfrac{12}{30},\ \dfrac{9}{30}$
  \task $\text{m.c.m.}=24$: \ $\dfrac{15}{24},\ \dfrac{12}{24},\ \dfrac{14}{24}$
  \task $\text{m.c.m.}=42$: \ $\dfrac{24}{42},\ \dfrac{14}{42},\ \dfrac{15}{42}$
  \task $\text{m.c.m.}=180$: \ $\dfrac{30}{180},\ \dfrac{24}{180},\ \dfrac{35}{180}$
  \task $\text{m.c.m.}=28$: \ $\dfrac{24}{28},\ \dfrac{4}{28},\ \dfrac{12}{28}$
  \task $\text{m.c.m.}=60$: \ $\dfrac{100}{60},\ \dfrac{25}{60},\ \dfrac{45}{60},\ \dfrac{27}{60}$
\end{tasks}

\ex{2} $\res{\dfrac{30}{36},\ \dfrac{21}{36},\ \dfrac{16}{36}}$

\ex{3} $\dfrac{2}{3},\ \dfrac{2}{3},\ \dfrac{3}{4}$ \ $\to$ \ $\text{m.c.m.}(3,3,4)=12$: \ $\res{\dfrac{8}{12},\ \dfrac{8}{12},\ \dfrac{9}{12}}$.


% =====================================================
\seccio{7}{Comparar y ordenar fracciones}

\ex{1} \begin{tasks}(3)
  \task $\dfrac{4}{7}<\dfrac{5}{7}<\dfrac{6}{7}$
  \task $\dfrac{3}{9}<\dfrac{4}{9}<\dfrac{5}{9}$
  \task $\dfrac{2}{13}<\dfrac{3}{13}<\dfrac{4}{13}$
\end{tasks}\par (Hay muchas respuestas posibles.)

\ex{1} \begin{tasks}(4)
  \task $\dfrac{9}{5}\ \res{>}\ \dfrac{9}{8}$
  \task $\dfrac{6}{11}\ \res{<}\ \dfrac{10}{11}$
  \task $\dfrac{3}{7}\ \res{<}\ \dfrac{3}{4}$
  \task $\dfrac{4}{8}\ \res{>}\ \dfrac{2}{16}$
  \task $\dfrac{2}{5}\ \res{<}\ \dfrac{5}{6}$
  \task $\dfrac{1}{3}\ \res{<}\ \dfrac{3}{7}$
  \task $\dfrac{8}{9}\ \res{<}\ \dfrac{11}{12}$
  \task $\dfrac{3}{4}\ \res{=}\ \dfrac{24}{32}$
\end{tasks}

\ex{2} \begin{tasks}(3)
  \task $\dfrac{2}{7}<\dfrac{3}{5}<\dfrac{11}{2}$
  \task $\dfrac{4}{9}<\dfrac{2}{3}<\dfrac{5}{6}$
  \task $\dfrac{1}{10}<\dfrac{2}{6}<\dfrac{5}{4}$
  \task $\dfrac{23}{21}<\dfrac{14}{7}<\dfrac{18}{3}$
  \task $\dfrac{3}{24}<\dfrac{2}{12}<\dfrac{4}{15}$
  \task $\dfrac{3}{28}<\dfrac{12}{14}=\dfrac{6}{7}$
\end{tasks}

\ex{2} \begin{tasks}(2)
  \task $\dfrac{2}{8}<\dfrac{3}{10}<\dfrac{7}{20}<\dfrac{4}{5}$
  \task $\dfrac{4}{7}<\dfrac{2}{3}<\dfrac{5}{4}<\dfrac{5}{2}$
  \task $\dfrac{1}{4}<\dfrac{2}{6}<\dfrac{7}{12}<\dfrac{3}{2}$
\end{tasks}

\ex{2} {\centering \begin{tikzpicture}[baseline=0]\draw[-latex,thick] (-0.3,0)--(9.50,0);\draw[muted] (1.125,0.1)--(1.125,-0.1);\draw[muted] (2.250,0.1)--(2.250,-0.1);\draw[muted] (3.375,0.1)--(3.375,-0.1);\draw[muted] (5.625,0.1)--(5.625,-0.1);\draw[muted] (6.750,0.1)--(6.750,-0.1);\draw[muted] (7.875,0.1)--(7.875,-0.1);\draw[thick] (0.00,0.18)--(0.00,-0.18) node[below=1pt,font=\small]{$0$};\draw[thick] (4.50,0.18)--(4.50,-0.18) node[below=1pt,font=\small]{$1$};\draw[thick] (9.00,0.18)--(9.00,-0.18) node[below=1pt,font=\small]{$2$};\fill[accent] (1.125,0) circle (0.09) node[above=3pt,text=accent]{$\frac{1}{4}$};\fill[accent] (3.375,0) circle (0.09) node[below=10pt=3pt,text=accent]{$\frac{3}{4}$};\fill[accent] (5.625,0) circle (0.09) node[above=3pt,text=accent]{$\frac{5}{4}$};\fill[accent] (6.750,0) circle (0.09) node[below=10pt=3pt,text=accent]{$\frac{3}{2}$};\end{tikzpicture}\par}

\ex{3} $\dfrac{2}{5}=\dfrac{4}{10}$ y $\dfrac{3}{5}=\dfrac{6}{10}$: \ $\res{\dfrac{5}{10}}$. \qquad $\dfrac{1}{3}=\dfrac{4}{12}$ y $\dfrac{1}{2}=\dfrac{6}{12}$: \ $\res{\dfrac{5}{12}}$.

\ex{3} $\text{m.c.m.}=40$: \ $\dfrac{30}{40},\ \dfrac{25}{40},\ \dfrac{28}{40}$. \ Ha leído más \res{\text{Ana}}, después Clara y después Gabriel.


% =====================================================
\seccio{8}{Sumas y restas}

\ex{1} \begin{tasks}(3)
  \task $\dfrac{4}{7}+\dfrac{8}{7}=\res{\dfrac{12}{7}}$
  \task $\dfrac{2}{3}+\dfrac{1}{5}=\res{\dfrac{13}{15}}$
  \task $\dfrac{5}{2}-\dfrac{1}{8}=\res{\dfrac{19}{8}}$
  \task $\dfrac{3}{6}-\dfrac{4}{12}=\res{\dfrac{1}{6}}$
  \task $\dfrac{5}{8}+\dfrac{1}{5}=\res{\dfrac{33}{40}}$
  \task $\dfrac{14}{15}-\dfrac{2}{3}=\res{\dfrac{4}{15}}$
\end{tasks}

\ex{1} \begin{tasks}(3)
  \task $\dfrac{3}{4}+\dfrac{6}{8}=\res{\dfrac{3}{2}}$
  \task $4+\dfrac{25}{9}=\res{\dfrac{61}{9}}$
  \task $\dfrac{25}{8}-2=\res{\dfrac{9}{8}}$
  \task $\dfrac{3}{5}+1=\res{\dfrac{8}{5}}$
  \task $5-\dfrac{3}{25}=\res{\dfrac{122}{25}}$
  \task $\dfrac{13}{6}+\dfrac{2}{7}=\res{\dfrac{103}{42}}$
\end{tasks}

\ex{2} \begin{tasks}(3)
  \task $\dfrac{12}{15}+\dfrac{8}{9}=\res{\dfrac{76}{45}}$
  \task $\dfrac{7}{12}-\dfrac{3}{10}=\res{\dfrac{17}{60}}$
  \task $\dfrac{7}{6}+\dfrac{8}{12}=\res{\dfrac{11}{6}}$
  \task $\dfrac{18}{5}-\dfrac{3}{10}=\res{\dfrac{33}{10}}$
  \task $\dfrac{5}{3}-\dfrac{9}{24}=\res{\dfrac{31}{24}}$
  \task $\dfrac{35}{36}+\dfrac{2}{9}=\res{\dfrac{43}{36}}$
\end{tasks}

\ex{2} \begin{tasks}(2)
  \task $\dfrac{4}{5}+\dfrac{2}{7}-\dfrac{1}{2}=\res{\dfrac{41}{70}}$
  \task $\dfrac{6}{8}+\dfrac{4}{3}-\dfrac{8}{12}=\res{\dfrac{17}{12}}$
  \task $\dfrac{3}{2}+\dfrac{6}{7}-\dfrac{2}{4}=\res{\dfrac{13}{7}}$
  \task $6-\dfrac{1}{2}-\dfrac{4}{12}=\res{\dfrac{31}{6}}$
  \task $\dfrac{2}{3}+\dfrac{1}{5}-\dfrac{1}{6}=\res{\dfrac{7}{10}}$
  \task $\dfrac{4}{3}+\dfrac{1}{6}-\dfrac{2}{7}=\res{\dfrac{17}{14}}$
\end{tasks}

\ex{3} \begin{tasks}(2)
  \task $\dfrac{2}{10}+\dfrac{5}{60}-\dfrac{1}{30}=\res{\dfrac{1}{4}}$
  \task $\dfrac{12}{5}+2-\dfrac{1}{25}=\res{\dfrac{109}{25}}$
  \task $\dfrac{6}{9}+\dfrac{1}{3}-\dfrac{2}{6}=\res{\dfrac{2}{3}}$
  \task $\dfrac{10}{3}-\dfrac{2}{6}-\dfrac{4}{5}+\dfrac{1}{2}=\res{\dfrac{27}{10}}$
  \task $4-\dfrac{1}{2}+\dfrac{1}{3}+2=\res{\dfrac{35}{6}}$
  \task $\dfrac{7}{9}+\dfrac{1}{3}-\dfrac{2}{5}=\res{\dfrac{32}{45}}$
\end{tasks}

\ex{3} $\dfrac{1}{4}+\dfrac{2}{5}+\dfrac{1}{6}=\dfrac{49}{60}$, \ falta $1-\dfrac{49}{60}=\res{\dfrac{11}{60}}$.


% =====================================================
\seccio{9}{Multiplicación y división}

\ex{1} \begin{tasks}(3)
  \task $\dfrac{6}{4}\cdot \dfrac{2}{5}=\res{\dfrac{3}{5}}$
  \task $\dfrac{3}{2}\cdot \dfrac{2}{3}=\res{1}$
  \task $2\cdot \dfrac{5}{6}=\res{\dfrac{5}{3}}$
  \task $\dfrac{3}{2}\cdot \dfrac{4}{7}\cdot 3=\res{\dfrac{18}{7}}$
  \task $\dfrac{3}{4}\cdot \dfrac{10}{2}\cdot \dfrac{1}{6}=\res{\dfrac{5}{8}}$
  \task $\dfrac{4}{3}\cdot \dfrac{2}{8}\cdot \dfrac{6}{2}=\res{1}$
\end{tasks}\par En el b) y el f) el producto es 1: las fracciones son \textbf{inversas}.

\ex{1} \begin{tasks}(4)
  \task $\res{\dfrac{2}{5}}$
  \task $\res{\dfrac{1}{5}}$
  \task $\res{4}$
  \task $\res{\dfrac{3}{7}}$
\end{tasks}

\ex{1} \begin{tasks}(3)
  \task $\dfrac{2}{5}\cdot\dfrac{5}{3}=\dfrac{2}{3}$: \ \res{\text{no}}
  \task $\dfrac{4}{3}\cdot\dfrac{12}{16}=1$: \ \res{\text{sí}}
  \task $\dfrac{7}{9}\cdot\dfrac{9}{7}=1$: \ \res{\text{sí}}
  \task $\dfrac{2}{4}\cdot\dfrac{4}{2}=1$: \ \res{\text{sí}}
  \task $\dfrac{5}{6}\cdot\dfrac{36}{30}=1$: \ \res{\text{sí}}
  \task $\dfrac{12}{13}\cdot\dfrac{26}{12}=2$: \ \res{\text{no}}
\end{tasks}

\ex{1} \begin{tasks}(4)
  \task $\dfrac{13}{4}$
  \task $\dfrac{7}{3}$
  \task $\dfrac{8}{15}$
  \task $\dfrac{1}{6}$
\end{tasks}

\ex{2} \begin{tasks}(3)
  \task $\dfrac{4}{3}:\dfrac{5}{6}=\res{\dfrac{8}{5}}$
  \task $\dfrac{5}{4}:2=\res{\dfrac{5}{8}}$
  \task $\dfrac{4}{8}:\dfrac{6}{3}=\res{\dfrac{1}{4}}$
  \task $\dfrac{2}{6}:\dfrac{5}{4}:\dfrac{1}{3}=\res{\dfrac{4}{5}}$
  \task $\dfrac{3}{4}:\dfrac{12}{5}:\dfrac{2}{5}=\res{\dfrac{25}{32}}$
  \task $6:\dfrac{2}{3}:\dfrac{1}{5}=\res{45}$
\end{tasks}

\ex{2} \begin{tasks}(2)
  \task $\dfrac{3}{4}\cdot \dfrac{5}{8}:\dfrac{2}{5}=\res{\dfrac{75}{64}}$
  \task $\dfrac{5}{4}:\dfrac{1}{3}\cdot \dfrac{6}{5}=\res{\dfrac{9}{2}}$
  \task $\dfrac{5}{12}\cdot \dfrac{3}{2}:\dfrac{1}{6}=\res{\dfrac{15}{4}}$
  \task $\dfrac{6}{5}:\dfrac{7}{4}\cdot \dfrac{2}{10}=\res{\dfrac{24}{175}}$
\end{tasks}

\ex{3} $12:\dfrac{3}{4}=\dfrac{48}{3}=\res{16}$ botellas. \qquad $\dfrac{9}{10}:\dfrac{3}{4}=\dfrac{36}{30}=\res{\dfrac{6}{5}}$.


% =====================================================
\seccio{10}{Operaciones combinadas}

\ex{2} \begin{tasks}(2)
  \task $\dfrac{5}{2}+\dfrac{1}{3}\cdot \dfrac{4}{5}=\res{\dfrac{83}{30}}$
  \task $\dfrac{21}{4}-\dfrac{2}{5}:\dfrac{4}{3}=\res{\dfrac{99}{20}}$
  \task $\dfrac{4}{3}:\dfrac{2}{6}+\dfrac{1}{3}\cdot \dfrac{5}{4}=\res{\dfrac{53}{12}}$
  \task $\dfrac{4}{3}\cdot \left(\dfrac{2}{5}+\dfrac{1}{3}\right)=\res{\dfrac{44}{45}}$
  \task $\left(\dfrac{2}{5}+1\right):\dfrac{8}{3}=\res{\dfrac{21}{40}}$
  \task $\left(\dfrac{5}{4}+1\right)\cdot \dfrac{6}{7}=\res{\dfrac{27}{14}}$
\end{tasks}

\ex{2} \begin{tasks}(2)
  \task $\left(\dfrac{2}{5}+\dfrac{3}{7}\right)\cdot \dfrac{1}{2}-\dfrac{4}{10}=\res{\dfrac{1}{70}}$
  \task $\left(\dfrac{5}{4}+\dfrac{2}{3}+\dfrac{1}{12}\right)\cdot \dfrac{5}{4}-1=\res{\dfrac{3}{2}}$
  \task $\dfrac{5}{4}:\left(\dfrac{6}{4}-1\right)+\dfrac{2}{5}=\res{\dfrac{29}{10}}$
  \task $\dfrac{4}{3}-\left(\dfrac{1}{2}+\dfrac{2}{7}\right)+\dfrac{4}{5}\cdot \dfrac{2}{4}=\res{\dfrac{199}{210}}$
\end{tasks}

\ex{3} \begin{tasks}(2)
  \task $\left(\dfrac{1}{4}+\dfrac{2}{3}-\dfrac{1}{12}\right)\cdot 3+\left(\dfrac{3}{2}-\dfrac{5}{4}\right):\dfrac{1}{2}=\res{3}$
  \task $4\cdot \left(\dfrac{2}{3}+\dfrac{4}{6}-1\right)+\dfrac{2}{3}=\res{2}$
  \task $\left(\dfrac{3}{2}+\dfrac{1}{7}\right)\cdot \left(\dfrac{5}{3}-\dfrac{7}{15}\right)=\res{\dfrac{69}{35}}$
  \task $\dfrac{12}{10}:\dfrac{5}{4}+\dfrac{1}{2}\cdot \dfrac{13}{5}=\res{\dfrac{113}{50}}$
  \task $\left(\dfrac{4}{3}+\dfrac{2}{5}\cdot \dfrac{1}{3}\right):\dfrac{4}{12}=\res{\dfrac{22}{5}}$
  \task $\dfrac{3}{2}:\dfrac{1}{4}\cdot \dfrac{4}{20}+\dfrac{3}{2}\cdot \dfrac{7}{5}=\res{\dfrac{33}{10}}$
\end{tasks}

\ex{3} \begin{tasks}(1)
  \task $\left(\dfrac{5}{4}+\dfrac{2}{3}-1\right)\cdot \dfrac{3}{5}+\left(\dfrac{2}{5}+\dfrac{4}{3}\right)\cdot 2=\res{\dfrac{241}{60}}$
  \task $\dfrac{3}{8}+\dfrac{4}{5}\cdot \dfrac{1}{3}+\left(\dfrac{2}{6}-\dfrac{1}{4}\right)\cdot 3=\res{\dfrac{107}{120}}$
\end{tasks}


% =====================================================
\seccio{11}{Problemas con fracciones}

\ex{1} \begin{tasks}(1)
  \task $\text{m.c.m.}=280$: \ $\dfrac{40}{280}+\dfrac{112}{280}+\dfrac{35}{280}=\res{\dfrac{187}{280}}$
  \task $1-\dfrac{187}{280}=\res{\dfrac{93}{280}}$
  \task Más: \res{\text{Elena}} ($\dfrac{112}{280}$); \ menos: \res{\text{Maite}} ($\dfrac{35}{280}$).
\end{tasks}

\ex{1} Nos hemos comido $(24:8)\cdot5=15$: \ quedan $\res{9}$ porciones.

\ex{2} \begin{tasks}(1)
  \task $\dfrac{1}{2}+\dfrac{1}{7}+\dfrac{2}{7}=\dfrac{13}{14}$: \ sobra $\res{\dfrac{1}{14}}$.
  \task $\dfrac{1}{14}+\dfrac{3}{5}=\dfrac{5}{70}+\dfrac{42}{70}=\res{\dfrac{47}{70}}$ de pizza.
\end{tasks}

\ex{2} \begin{tasks}(1)
  \task $1-\dfrac{1}{2}-\dfrac{1}{5}=\res{\dfrac{3}{10}}$
  \task Más de 65: $\res{2625}$; \ menos de 18: $\res{1050}$; \ entre 18 y 65: $\res{1575}$.
\end{tasks}

\ex{2} Se gastan 30 L (quedan 30); \ después $\dfrac{1}{3}$ de 30 $=10$ L (quedan 20); \ se derrama la mitad: quedan $\res{10\text{ L}}$.

\ex{3} $\dfrac{35}{105}+\dfrac{21}{105}+\dfrac{30}{105}=\dfrac{86}{105}$ \ $\Rightarrow$ \ $1-\dfrac{86}{105}=\res{\dfrac{19}{105}}$.

\ex{3} \begin{tasks}(1)
  \task Hipoteca $\res{1800\text{\,€}}$, comida $\res{1080\text{\,€}}$, coche $\res{600\text{\,€}}$; \ quedan $\res{1920\text{\,€}}$.
  \task $\dfrac{1920}{5400}=\res{\dfrac{16}{45}}$
\end{tasks}

\ex{3} Queda $\dfrac{1}{3}$ de $\dfrac{3}{4}$, es decir, $\dfrac{1}{4}$ del depósito, que son 50 L: \ capacidad $50\cdot4=\res{200\text{ L}}$.


\end{document}
